\chapter{\texorpdfstring{$\Aut_E$}{Aut（E）} 的构造}\label{ch:8}
\srcpage{44}
\section{已经完成的构造}\label{sec:8.1}
继续取 $G=\GL_n$，并令 $E$ 为 $X$ 上几何不可约的秩 $n$ 局部系统。第\ref{ch:5}–\ref{ch:7}章在良好的开子叠 $C'_n\subset\Coh'_n$ 上构造了初步的 Whittaker–Laumon 层
\[\Aut'_E\in D(C'_n).\]
它通过公式
\[\Aut'_E=\pi_!(p_1^*\cL_E\otimes p_2^*\cL_\psi)[\text{移位}](\text{扭曲})\]
从紧化 Whittaker 叠
\origdiagram{p044-diagram.pdf}{0.4258}
得到。已经证明的定理是：在 $C'_n$ 上，这个对象是不可约反常层。

令
\[C_n^\circ:=C_n\cap\Bun_n,\qquad C_n^{\prime,\circ}:=C'_n\cap\Bun'_n,\]
并令
\[p_C:C_n^{\prime,\circ}\longrightarrow C_n^\circ\]
为将带有一般点非零截面 $\Omega_X^{n-1}\to M$ 的丛映为丛 $M$ 的遗忘映射。待证明的下降陈述是：若 $C_n\cap\Bun_n^d\ne\varnothing$，则
\[\Aut'_E|_{C_n^{\prime,\circ}\cap(\Bun_n^d)'}\]
沿 $p_C$ 下降为层
\[\Aut_E|_{C_n\cap\Bun_n^d}\in\operatorname{Perv}(C_n\cap\Bun_n^d).\]
由于 $p_C$ 在所选良好开子叠上光滑，非零的下降层自动是不可约反常层。最终目标是将这些局部部分唯一延拓为整个 $\Bun_n$ 上的 Hecke 特征层 $\Aut_E$。
\section{沿射影丛纤维下降}\label{sec:8.2}\index{下降!射影丛|hyperpage}
在光滑开图 $V\subset C_n^\circ\cap\Bun_n^d$ 上工作。必要时缩小 $V$ 后，截面 $\Omega_X^{n-1}\to M$ 的空间是 $V$ 上的向量丛；再对截面上的标量作用取商，得到射影丛型映射
\[p_V:V'\longrightarrow V.\]
下降的图景是：先将 $\Aut'_E|_{V'}$ 写为 $V'$ 的某个稠密开部分上局部系统的 IC 延拓，再证明这个 IC 延拓实际上是整个 $V'$ 上的局部系统。

使用两个条件。首先，对每个几何点 $m\in V$，良好开部分与纤维的交是稠密的：
\[V'\cap p_V^{-1}(m)\subset p_V^{-1}(m)\quad\text{是开且稠密的。}\]
可以将它纳入良好开部分的定义。其次，$\Aut'_E|_{V'}$ 必须是局部系统，而不只是 IC 层。知道这一点后，局部系统沿 $p_V$ 唯一下降，因为几何纤维是射影空间，且
\[\pi_1(\mathbb P_{\bar k}^m)=1.\]
\srcpage{45}
等价地，从 $V$ 上局部系统到 $V'$ 上局部系统的拉回函子是等价。这给出了所需的从 $C_n^{\prime,\circ}$ 到 $C_n^\circ$ 的下降。

用于证明不可约反常层为局部系统的判据，是当前纯 IC 情形中的有限域迹检验。设 $Y$ 光滑，$\cK\in\operatorname{Perv}(Y)$ 不可约。记
\[\chi_{\cK}(y):=\sum_i(-1)^i\dim H^i(\cK_{\bar y})\]
为茎的 Euler 示性数。对任意反常层，这个函数本身并不是局部系统判据。这里 $\cK$ 是由该构造得到的不可约 IC 延拓；在这一情形，$\chi_{\cK}$ 恒定且非零排除了边界贡献，从而迫使 $\cK$ 是 $Y$ 上移位的局部系统。因此，只需证明：
\[\chi(\Aut'_E)\text{ 在 }\Bun'_n\to\Bun_n\text{ 的纤维上恒定},\qquad\chi(\Aut'_E)\ne0\text{ 在 }C_n^\circ\text{ 上。}\]
\section{沿纤维的迹恒定性}\label{sec:8.3}
$\chi(\Aut'_E)$ 的恒定性归约到已知存在自守函数的情形。作为黑箱输入，使用 Deligne 的真推前不变性：若 $f:Y_1\to Y_2$ 为真态射，$\cK_1,\cK_2\in D(Y_1)$ 在 étale 局部同构，则 $f_!\cK_1$ 与 $f_!\cK_2$ 的 Euler 示性数函数相同。
\begin{transnote}{含糊点拨}
Euler 示性数函数可按分层计算：把 $Y_1$ 细分到 $\cK_1$ 与 $\cK_2$ 都局部常值（带平移与扭转）的层上，两层在每个分层上只差一个常数因子，故 $\chi$ 相同；真性保证 $f_!\cK_i$ 仍是构造性的，$\chi$ 才有定义。书中把它当黑箱，是因为只需借它把「沿 $\Bun'_n\to\Bun_n$ 纤维恒定」转移到某一个方便的自守函数上，并不需要 Deligne 定理的内部机制。[来源：nLab "six operations"，https://ncatlab.org/nlab/show/six+operations（Last revised 2026-09-30；2026-11-17 访问）；并见原书 p.45 本段。]
\end{transnote}

对任意两个秩 $n$ 局部系统 $E_1,E_2$，即使它们可约，在 $X$ 上也 étale 局部同构。因此，相应的 Laumon 层 étale 局部同构：
\[\cL_{E_1}\simeq_{\mathrm{et,loc}}\cL_{E_2}.\]
若定义 $\Aut'_E$ 的映射是真态射，Deligne 定理立即给出
\[\chi(\Aut'_{E_1})=\chi(\Aut'_{E_2}).\]
该映射并非字面上的真态射，但紧化 Whittaker 构造有 Drinfeld 紧化，所以将对应替换为其紧化版本后，同样的论证适用。

因此，对给定 $E$ 证明恒定性命题，只需对方便的秩 $n$ 局部系统 $E'$ 证明。选择 $E'$，使存在以 $E'$ 为 Hecke 特征值的尖点 Hecke 特征函数。由 Frenkel–Gaitsgory–Kazhdan–Vilonen 的比较定理，它在 $\Bun'_n(\mathbb F_q)$ 上的限制是迹函数
\[M'\longmapsto\Tr(\operatorname{Frob}_{M'},\Aut'_{E'}).\]
\begin{transnote}{含糊点拨}
比较定理的两边是：几何侧的特征层 $\Aut'_{E'}$ 在 $\Bun'_n(\mathbb F_q)$ 上的迹函数，与自守侧的 Whittaker 归一化特征函数。定理说这两者相等（在 $E'$ 为 Hecke 特征值的前提下），于是「数点」与「层」可以互换。这里限制到 $\Bun'_n$ 而不是 $\Bun_n$，是因为 Whittaker 归一化正落在 $\Bun'_n$ 上；$E'$ 之所以要选成「以尖点 Hecke 特征函数为特征值」，是为了让迹函数非零且可算。[依据：原书 p.45 本段；并见 [FGKV98]（JAMS 11 (1998), 451–484）的层–函数比较定理。]
\end{transnote}
由于这个函数从 $\Bun_n(\mathbb F_q)$ 拉回，它在 $\Bun'_n\to\Bun_n$ 的每个纤维上恒定。由前述归约，同样的恒定性对 $\Aut'_E$ 成立。这证明了下降定理。
\section{Hecke–Laumon 对应}\label{sec:8.4}\index{Hecke–Laumon 对应|hyperpage}
通常的 Hecke 性质由凝聚层上一个看起来更强的性质推出。对 $k\ge0$、$d\ge0$，定义 Hecke–Laumon 叠
\[\mathrm{HL}_k^d:=\{0\to M'\to M\to\mathcal T\to0\mid M,M'\in\Coh_k,\ \mathcal T\in\Tor^d\}.\]
它具有映射
\[\Coh_k\xleftarrow{\overleftarrow h}\mathrm{HL}_k^d\xrightarrow{\overrightarrow h}\Tor^d\times\Coh_k,\]
其中 $\overleftarrow h$ 记住 $M$，$\overrightarrow h$ 记住 $(\mathcal T,M')$。相应函子为
\[\mathrm{HL}_k^d(\cK):=\overrightarrow h_!\overleftarrow h^*\cK[dk](dk/2),\qquad\cK\in D(\Coh_k).\]
\begin{transnote}{易错预警}
两个因子各管一件事，缺一个就错：$[dk]$ 是 $\mathrm{HL}_k^d$ 在 $\Coh_k$ 上的相对维数（每次修改贡献 $k$，共 $d$ 次），$(dk/2)$ 是配对的 Tate 归一化，与 §5.2 的 $[r](r/2)$ 同一个来源。两者合起来才使 $\mathrm{HL}_k^d$ 对反常 $t$-结构 $t$-正合；$(dk/2)$ 可以是半整数，故系数域须含 $\sqrt q$。[来源：nLab "Tate twist"，https://ncatlab.org/nlab/show/Tate+twist（Last revised 2025-08-10；2026-11-17 访问）；并见原书 p.45–46 本式与 §5.2 的归一化约定。]
\end{transnote}
\srcpage{46}
\begin{definition*}{}
若反常层 $\cK\in\operatorname{Perv}(\Coh_k)$ 配备同构
\[\alpha_d:\mathrm{HL}_k^d(\cK)\xrightarrow{\sim}\cL_E^d\boxtimes\cK\qquad(d\ge0),\]
满足 $d=0$ 时显然的归一化，且这些同构与迭代修改相容，则称它为以 $E$ 为特征值的 Hecke–Laumon 特征层。
\end{definition*}
更明确地，对 $d=d_1+d_2$，比较通过
\[0\to M'\to M\to\mathcal T'\to0,\qquad0\to M''\to M'\to\mathcal T''\to0\]
从 $M$ 转到 $M''$ 的两种方式。这种两步修改的叠既映到 $\mathrm{HL}_k^d$，也映到 $\Tor^{d_1}\times\mathrm{HL}_k^{d_2}$。所需相容性要求由
\[\mathrm{HL}_k^{d_2}\circ\mathrm{HL}_k^{d_1}\simeq((s,q)_!m^*\boxtimes\mathrm{Id})\circ\mathrm{HL}_k^d\]
得到的图与\chref{7}的因子化同构
\[(s,q)_!m^*\cL_E^d\simeq\cL_E^{d_1}\boxtimes\cL_E^{d_2}\]
交换。这里 $m,s,q$ 是挠层短正合列叠
\[0\longrightarrow\mathcal T''\longrightarrow\mathcal T\longrightarrow\mathcal T'\longrightarrow0\]
中的映射。

对 $k=0$，这恰好是 Laumon 因子化定理：$\cL_E\in\operatorname{Perv}(\Tor)$ 是 Hecke–Laumon 特征层。前几章的 Fourier 归纳继而对 $k$ 归纳，给出
\[F_k\in\operatorname{Perv}(C'_k)\]
为 Hecke–Laumon 特征层。特别地，$\Aut'_E|_{C_n^{\prime,\circ}}$ 满足 Hecke–Laumon 特征性质。下降后，层 $\Aut_E|_{C_n^\circ}$ 继承同样的性质。
\section{从 Hecke–Laumon 到通常的 Hecke}\label{sec:8.5}
剩下的问题是，为何 Hecke–Laumon 性质蕴涵 $\Bun_n$ 上通常的 Hecke 特征性质。这里使用的陈述为：
\begin{proposition}{}{8.5.1}
设 $\mathcal S\in\operatorname{Perv}(C_n^\circ)$ 是以 $E$ 为特征值的 Hecke–Laumon 特征层。再假设其 Verdier 对偶 $\mathbb D\mathcal S$ 是以 $E^*$ 为特征值的 Hecke–Laumon 特征层。则 $\mathcal S$ 满足以 $E$ 为特征值的通常 Hecke 特征性质。
\end{proposition}
\begin{proof}

只需证明第一 Hecke 关系。考虑 $\mathrm{HL}_n^1$。去掉分裂的短正合列，得到开子叠
\[\mathrm{HL}_n^{1,\circ}:=\{0\to M'\to M\to\mathcal T\to0\mid\length(\mathcal T)=1,\ \text{序列不分裂}\}.\]
沿 $X\to\Tor^1$ 作基变换；该映射将 $x$ 映为支撑于 $x$ 的长度一挠层。所得叠为通常的第一 Hecke 对应：
\[\mathrm{HL}_n^{1,\circ}\times_{\Tor^1}X=\{(x,0\to M'\to M)\mid x\in X,\ \supp(M/M')=x,\ \length(M/M')=1\}.\]
等价地，
\[M'\subset M\subset M'(x),\qquad\length(M/M')=1.\]
长度一挠商剩余的标量自同构是一个 $\Gm$ 因子；对它取商后得到通常的 Hecke 叠 $\Hecke_n^1$。

这一商所需的反常估计由下面的初等引理提供。
\end{proof}
\srcpage{47}
\begin{lemma}{}{8.5.2}
设 $\mathcal E\xrightarrow{\pi}B$ 为向量丛，$\mathbb P(\mathcal E)\xrightarrow{\bar\pi}B$ 为其射影化。假设 $\cK\in\operatorname{Perv}(\mathcal E)$ 是 $\Gm$-等变的，且 $\pi_!\cK,\pi_*\cK\in\operatorname{Perv}(B)$。令 $\overline{\cK}$ 为 $\cK|_{\mathcal E\setminus0}$ 到 $\mathbb P(\mathcal E)$ 的下降。则 $\bar\pi_!\overline{\cK}$ 是反常层。
\end{lemma}
\begin{proof}

将\lemref{8.5.2}逐纤维应用于从 $\mathrm{HL}_n^{1,\circ}$ 到通常 Hecke 叠的商。\secref{8.4}中的 Hecke–Laumon 特征性质给出 $\pi_!\cK$ 所需的反常性，对 $\mathbb D\mathcal S$ 的同一性质给出 $\pi_*\cK$ 的反常性。由于
\[\Tor^1\simeq X\times B\Gm\]
且 $\cL_E^1|_X=E$，$\mathcal S$ 的第一通常 Hecke 变换是反常的，且特征值为 $E$。较高的微小 Hecke 关系由第一 Hecke 对应在卷积下的相容性得到，特征值为 $\bigwedge^iE$。因此，$\mathcal S$ 是通常的 Hecke 特征层。
\end{proof}

对 $\mathcal S=\Aut_E|_{C_n^\circ}$，Hecke–Laumon 归纳与 Verdier 对偶保证这些假设成立：$\mathbb D(\Aut_E)$ 的特征值为 $E^*$。所以 $\Aut_E|_{C_n^\circ}$ 满足通常的 Hecke 特征性质。特别地，采用\chref{2}的 Hecke 对应约定，
\[(t_x^{d+n})^*\Aut_E^d\simeq\det(E_x)\otimes\Aut_E^{d+n},\]
其中 $t_x^{d+n}:\Bun_n^{d+n}\to\Bun_n^d$ 为同构，$M\mapsto M(-x)$。等价地，对 $\tau_x^d:\Bun_n^d\to\Bun_n^{d+n}$，$M\mapsto M(x)$，有
\[(\tau_x^d)^*\Aut_E^{d+n}\simeq\det(E_x)^{-1}\otimes\Aut_E^d.\]
\newpage
\section{延拓到整个 \texorpdfstring{$\Bun_n$}{Bunₙ} 及尖点性}\label{sec:8.6}
令 $j_d:C_n\cap\Bun_n^d\hookrightarrow\Bun_n^d$ 为次数 $d$ 分支中的良好开子叠。构造首先通过中间延拓扩展下降层：
\[\Aut_E^d:=j_{d,!*}(\Aut_E|_{C_n\cap\Bun_n^d}).\]
在洁净范围内，它与 $j_{d,!}$ 和 $j_{d,*}$ 给出同一对象。然后，上述最高次 Hecke 关系给出次数分支之间的唯一延拓，因为 $M\mapsto M(-x)$ 使次数减少 $n$。$\Hecke_n^1$ 与最高次 Hecke 对应之间的相容性，将第一 Hecke 关系从良好开子叠传播到延拓层。于是得到以 $E$ 为特征值的 Hecke 特征层
\[\Aut_E\in D(\Bun_n).\]

最后，$\Aut_E$ 是尖点层。原因是洁净延拓论证中所用的同一个消失定理：$\GL_n$ 的每个真抛物常数项函子，都经由某个 $k<n$ 的 $\Bun_k$ 上的平均操作分解，而几何不可约秩 $n$ 局部系统 $E$ 的平均定理迫使该贡献消失。因此，$\Aut_E$ 的所有真抛物常数项均消失。
\begin{transnote}{含糊点拨}
「真抛物常数项函子」是沿抛物子群丛 $\Bun_P\to\Bun_n$ 的 $!$-推前再拉回的复合，对应经典自守表示里的常数项。对 $\GL_n$，一个真抛物 $P$ 的 Levi 是 $\GL_{d_1}\times\cdots\times\GL_{d_r}$，其中必有一块 $k<n$，所以常数项可以分解成先到某个 $\Bun_k$、再作平均的形式；$\rank E=n$ 且几何不可约时，平均定理（定理 5.7.2）正给出这一项为零。这就是「$\Aut_E$ 尖点」的全部内容：不是逐项硬算，而是把每个真抛物常数项化成同一个消失定理。[依据：原书 p.47 本段；并见定理 5.7.2 与 [Gai16b]。]
\end{transnote}

在有限域上，与 FGKV 的层–函数比较陈述如下。若给定 $\Bun_n(\mathbb F_q)$ 上以 $E$ 为特征值的尖点 Hecke 特征函数，则它在 $\Bun'_n(\mathbb F_q)$ 上的限制为
\[M'\longmapsto\Tr(\operatorname{Frob}_{M'},\Aut'_E).\]
因此，上文构造的层 $\Aut_E$ 将预期自守特征函数几何化，与文献\cite{FGKV98}中 Whittaker 函数的比较相匹配。
